\documentclass[12pt]{article}
\usepackage[utf8]{inputenc}
\usepackage[T1]{fontenc}
\usepackage[margin=1in]{geometry}
\usepackage{setspace}
\usepackage{xcolor}
\usepackage[colorlinks=true, urlcolor=blue!60!black, linkcolor=blue!60!black]{hyperref}
\setlength{\parindent}{0em}
\setlength{\parskip}{0.6em}
\renewcommand{\baselinestretch}{1.0}
\begin{document}
\begin{center}
{\LARGE \textit{Geduldig v. Aiello}, 417 U.S. 484 (1974)\par}
\vspace{0.5em}
{\large Case Brief --- Two-Page Edition\par}
\end{center}

\section{Memory Jogger}
Excluding normal pregnancy from a state disability insurance program is not sex discrimination: it divides ``pregnant women and nonpregnant persons'' (\href{https://supreme.justia.com/cases/federal/us/417/484/}{Justia}).

\section{Detailed Case Facts}
California's self-supporting, employee-funded disability insurance program excluded disabilities ``arising in connection with pregnancy.'' After the State began paying abnormal-pregnancy claims, only Jacqueline Jaramillo's normal-pregnancy claim remained.

\section{Procedural History}
A three-judge district court held the exclusion irrational under the Equal Protection Clause; the State appealed directly.

\section{Judicial Votes}
Reversed 6--3. Stewart wrote for Burger, White, Blackmun, Powell, and Rehnquist; Brennan dissented, joined by Douglas and Marshall.

\section{Holding}
A pregnancy exclusion is not a sex-based classification; absent proof that it is a mere pretext for invidious sex discrimination, it receives rational basis review, which California's fiscal interests satisfied.

\section{Analysis of Opinions}
\textbf{Majority:} Covering normal pregnancy would raise costs by roughly a third and destroy the fund's self-supporting design; a State need not attack every aspect of a problem at once. The program insured the same risks for both sexes, and pregnancy is a distinct physical condition, not a proxy for sex. \textbf{Dissent:} the fund covered nearly every other disability, including male-only conditions such as prostatectomies, so singling out pregnancy created a double standard that ``inevitably constitutes sex discrimination.''

\section{Examples: Future Applications}
\textbf{Same-Side:} \textit{Skrmetti} applied the same logic to a ban on puberty blockers for gender dysphoria. \textbf{Opposite-Side:} refusing to hire previously pregnant women while hiring fathers is a transparent pretext. \textbf{Fence-Sitter:} a health plan excluding contraceptives but covering erectile-dysfunction drugs.

\section{Critique}
Critics call the ``pregnant women and nonpregnant persons'' syllogism strained, since pregnancy is a perfect proxy for sex, and argue it treats men as the norm. Congress rejected the Title VII sequel, \textit{Gilbert}, in the Pregnancy Discrimination Act.

\section{Key Quotations}
\begin{quote}
``While it is true that only women can become pregnant, it does not follow that every legislative classification concerning pregnancy is a sex-based classification like those considered in \textit{Reed}, \textit{supra}, and \textit{Frontiero}, \textit{supra}. Normal pregnancy is an objectively identifiable physical condition with unique characteristics. Absent a showing that distinctions involving pregnancy are mere pretexts designed to effect an invidious discrimination against the members of one sex or the other, lawmakers are constitutionally free to include or exclude pregnancy from the coverage of legislation such as this on any reasonable basis, just as with respect to any other physical condition.''
\end{quote}
\noindent --- Justice Stewart, for the Court, 417 U.S. at 496--497 n. 20.

\end{document}
